\begin{table}[!htbp]
\centering\small
\caption{部分泛化边界的逐次末段合作范围；四次运行全部保留}
\label{tab:w2-parameter-boundaries}
\begin{tabular}{cclrr}
\toprule
$r$ & 异常比例 & 方法 & 末段均值 & 末段范围 \\
\midrule
3 & 0\% & 完整模型 & 0.00502 & 0.00500--0.00505 \\
3 & 0\% & 固定评分/学习门控 & 0.00519 & 0.00517--0.00523 \\
3 & 0\% & 固定合作规则 & 0.00567 & 0.00553--0.00603 \\
3.6 & 10\% & 完整模型 & 0.81877 & 0.81692--0.82111 \\
3.6 & 10\% & 固定评分/学习门控 & 0.86113 & 0.85947--0.86181 \\
3.6 & 10\% & 固定合作规则 & 0.86162 & 0.86074--0.86293 \\
4.8 & 50\% & 完整模型 & 0.93973 & 0.93824--0.94072 \\
4.8 & 50\% & 固定评分/学习门控 & 0.96896 & 0.96804--0.96953 \\
4.8 & 50\% & 固定合作规则 & 0.94281 & 0.94240--0.94346 \\
\bottomrule
\end{tabular}
\par\smallskip{\zihao{-5}\rmfamily 来源：作者，本文模型开发与仿真实验，2026年。\par}
\end{table}
